\documentclass[11pt,a4paper]{article}
\usepackage[T1]{fontenc}
\usepackage[utf8]{inputenc}
\usepackage[main=italian,provide=*]{babel}
\usepackage{amsmath,amssymb}
\usepackage{geometry}
\geometry{margin=2.5cm}
\usepackage{booktabs}
\usepackage{seqsplit}
\usepackage{hyperref}
\hypersetup{colorlinks=true,linkcolor=blue,citecolor=blue,urlcolor=blue}

\title{Benchmark esatto (Fase 2) — dimero}
\author{Note di lavoro — Tesi triennale, Samuele Schiavi\\ Relatore: Prof.\ Alessandro Chiesa}
\date{}

\begin{document}
\maketitle

\section*{Introduzione}

Questo documento raccoglie il modulo di benchmark esatto per il dimero di spin-1/2
($N=2$). È il modulo su cui tutto il resto del lavoro sul dimero (VQE, Trotter,
correlazioni dinamiche) si appoggia senza mai modificarlo — la diagonalizzazione
numerica (\texttt{eigh}) è il metro di misura con cui ogni risultato variazionale
viene giudicato.

Per ciascun file: cosa fa, con quali verifiche è stato controllato in questa sessione,
e cosa se ne può concludere — per il modulo \texttt{.py}, anche le funzioni interne e
come si usa il file.

\section*{\texttt{dimer\_exact.py}}

Modulo di benchmark: Hamiltoniana come \texttt{SparsePauliOp}
($H=J(X_1X_2+Y_1Y_2+Z_1Z_2)+b(Z_1+Z_2)+D(X_1Z_2-Z_1X_2)$), spettro analitico per il
caso isotropo $D=0$ (singoletto + tripletto $M=+1,0,-1$), sweep su griglia di campo,
self-test interno.

\textbf{Funzioni interne.}
\begin{itemize}
\item \texttt{dimer\_hamiltonian(b,\ J,\ D)} — l'Hamiltoniana completa come
\texttt{SparsePauliOp}, sette termini di Pauli (\texttt{XX}, \texttt{YY},
\texttt{ZZ}, \texttt{ZI}, \texttt{IZ}, \texttt{XZ}, \texttt{ZX}); sorgente unica
riusata da tutti i moduli VQE del dimero.
\item \texttt{magnetization\_operator()} — $M_z=(Z_1+Z_2)/2$, convenzione di spin
($s=\sigma/2$).
\item \texttt{analytic\_eigenvalues(b,\ J)} — le quattro energie esatte per $D=0$
(singoletto $E=-3J$; tripletto $E=J+2bM$, $M=+1,0,-1$).
\item \texttt{exact\_sweep(b\_values,\ J,\ D)} — diagonalizza $H$ su una griglia di
campo $b$, restituisce spettro completo e osservabili del fondamentale.
\item \texttt{\_self\_test(J)} — confronta spettro numerico e analitico ($D=0$) su
tutto il range di $b$, a precisione macchina.
\end{itemize}

\textbf{Come si usa.} Come libreria: \texttt{from dimer\_exact import
dimer\_hamiltonian, exact\_sweep, \ldots} — è la dipendenza comune a
\texttt{vqe\_dimer.py} e a tutti gli altri moduli VQE del dimero. Come script
standalone: \texttt{python3 dimer\_exact.py} esegue il self-test e rigenera la
figura \texttt{dimer\_exact.png}.

\textbf{Verifiche.} Rieseguito in questa sessione: self-test a precisione macchina
($\max|E_\text{num}-E_\text{analitico}|=8.88\times10^{-16}$), figura rigenerata.

\textbf{Conclusione.} Modulo maturo e stabile: nessuna decomposizione ulteriore
necessaria, lo spettro si ottiene direttamente dalla composizione di due spin-1/2 in
forma chiusa.

\section*{\texttt{dimer\_exact.png}}

Figura generata dal modulo precedente: a sinistra lo spettro completo (quattro
livelli) con il fondamentale evidenziato sia per $D=0$ (incrocio vero a $B/J=2$) sia
per $D=0.2$ (anticrossing, gap aperto); a destra la magnetizzazione del
fondamentale, con lo stesso confronto — salto discontinuo per $D=0$, transizione
liscia (crossover) per $D\neq0$.

\textbf{Considerazioni.} Il confronto diretto fra $D=0$ (vero incrocio, senza
degenerazione tranne che esattamente a $B/J=2$) e $D\neq0$ (anticrossing, gap sempre
aperto, nessuna degenerazione mai) è la manifestazione più semplice possibile della
regola di non-incrocio di von Neumann–Wigner: un solo parametro reale libero ($b$),
quindi l'incrocio si apre non appena si accende un termine ($D$) che collega i due
livelli.

\section*{\texttt{dimer\_exact\_spiegato.ipynb}}

Guida interattiva al modulo precedente: perché \texttt{SparsePauliOp} è la sorgente
unica di $H$ (con verifica esplicita della convenzione delle label,
\texttt{"ZI"}$=Z\otimes I$), perché si usa \texttt{eigh} e non \texttt{eig} (matrice
hermitiana $\Rightarrow$ autovalori reali e autovettori ortonormali garantiti,
verificato numericamente), la distinzione fra il $J$ del codice (coefficiente della
somma di Pauli $X_1X_2+Y_1Y_2+Z_1Z_2$) e il $J=4t^2/U$ della derivazione da Hubbard
(differiscono per un fattore 4), il resto del modulo (\texttt{exact\_sweep},
\texttt{\_self\_test}, la figura), slider interattivi per esplorare $b$, $D$, e il
crossover incrocio/anticrossing, un riepilogo finale in "tre risposte".

\textbf{Verifiche.} Eseguito per intero in questa sessione (30 celle, zero errori):
Hamiltoniana verificata hermitiana, convenzione delle label confermata
(\texttt{"ZI"==kron(Z,I)}), self-test a precisione macchina
($8.88\times10^{-16}$), tutti gli slider e le figure generate correttamente.

\textbf{Considerazioni.} Il notebook si sofferma esplicitamente su due questioni
facilmente date per scontate: la differenza fra \texttt{eig} e \texttt{eigh}
(rilevante ogni volta che si diagonalizza una matrice hermitiana), e la distinzione
fra le due convenzioni per $J$ (quella del codice, coefficiente della somma di Pauli,
e quella della derivazione fisica da Hubbard) — un punto di ambiguità facile da
introdurre involontariamente se non documentato esplicitamente.

\textbf{Conclusione.} Guida completa e autosufficiente al modulo di benchmark,
pensata per essere letta da chi non ha già familiarità con le scelte implementative
(perché quella label, perché quella routine di diagonalizzazione, perché quella
convenzione per $J$).

\section*{Conclusione generale}

Il benchmark esatto del dimero è un modulo maturo e stabile, verificato a precisione
macchina, che funge da sorgente unica e non modificabile dell'Hamiltoniana per tutto
il lavoro successivo (VQE, e in futuro Trotter e correlazioni dinamiche). Non
richiede alcuna decomposizione oltre alla diagonalizzazione diretta: lo spettro del
dimero si scrive in forma chiusa immediatamente dalla composizione di due spin-1/2.

\end{document}
